%%%%%%%%%%%%%%%%%%%%%%%%%%%%%%%%%%%%%%%%%%%%%%%%%%%%%%%%%%%%%%%%%%%%%%%%%%%%%%
%  QUESTIONNAIRE DE L'ÉTUDE QUANTITATIVE -- document autonome (A4)
%
%  Compilation : pdflatex questionnaire  (deux passes, aucune bibliographie).
%  Ce fichier est indépendant de la présentation (main.tex) ; il en reprend la numérotation
%  des questions (Q1 à Q23) et des sections (1 à 9) : voir la partie « Étude quantitative ».
%
%  Il contient :
%   - la notice pour l'équipe (objectif, cibles, tailles d'échantillon, règles de saisie) ;
%   - le questionnaire tel que le voit le répondant (texte d'accueil, 9 sections, concept neutre) ;
%   - la variante « secours et gestion de crise » (questions reformulées) ;
%   - les indicateurs de lecture et les seuils proposés.
%
%  Ce que vous pouvez modifier sans risque : les tailles d'échantillon, les seuils (annexe B),
%  les prix de la question 21, la liste des activités (Q4) et des lieux d'achat (Q22).
%  Les passages en italique gris entre crochets sont des consignes de programmation (logique de saut,
%  tirage au hasard) : ils ne s'affichent pas au répondant.
%%%%%%%%%%%%%%%%%%%%%%%%%%%%%%%%%%%%%%%%%%%%%%%%%%%%%%%%%%%%%%%%%%%%%%%%%%%%%%
\documentclass[11pt,a4paper]{article}

\usepackage[T1]{fontenc}
\usepackage[utf8]{inputenc}
\usepackage{lmodern}
\usepackage[french]{babel}
\usepackage[margin=2.2cm,top=2cm,bottom=2cm]{geometry}
\usepackage{amssymb}
\usepackage{xcolor}
\usepackage{booktabs}
\usepackage{array}
\usepackage{tabularx}
\usepackage{enumitem}
\usepackage{multicol}
\usepackage{needspace}
\usepackage{fancyhdr}
\usepackage[colorlinks=true,linkcolor=black,urlcolor=black]{hyperref}

\definecolor{insarouge}{HTML}{E42618}
\definecolor{insagris}{HTML}{5F5F5F}
\definecolor{insabeige}{HTML}{F8F0EC}

\setlength{\parindent}{0pt}
\setlength{\parskip}{4pt}
\pagestyle{fancy}
\fancyhf{}
\renewcommand{\headrulewidth}{0pt}
\fancyfoot[C]{\small\color{insagris}Étude quantitative~: communiquer sans réseau mobile\quad --\quad page \thepage}

% titres de sections : le numéro de section de la présentation est celui du questionnaire
\newcommand{\sect}[2]{\par\Needspace{9\baselineskip}\vspace{10pt}{\large\sffamily\bfseries\color{insarouge}Section #1\ \ #2}\par\vspace{2pt}\hrule height 0.6pt\vspace{4pt}}
% consigne de programmation (invisible pour le répondant)
\newcommand{\qa}[1]{\textbf{\textcolor{insarouge}{#1}}}
\newcommand{\logique}[1]{{\small\itshape\color{insagris}[#1]}}
% question : \quest{numéro}{énoncé}
\newcommand{\quest}[2]{\par\Needspace{6\baselineskip}\vspace{4pt}\noindent{\sffamily\bfseries\color{insarouge}Q#1}\hspace{0.5em}\textbf{#2}\par\nopagebreak}
% liste de cases à cocher
\newenvironment{choix}
  {\begin{list}{$\square$}{\setlength{\itemsep}{0pt}\setlength{\topsep}{2pt}\setlength{\parsep}{0pt}%
    \setlength{\leftmargin}{1.7em}\setlength{\labelwidth}{1.2em}\setlength{\labelsep}{0.5em}}}
  {\end{list}}
% même liste sur deux colonnes (pour les listes longues)
\newenvironment{choixdeux}
  {\begin{multicols}{2}\begin{choix}}
  {\end{choix}\end{multicols}}
% échelle à cinq niveaux : \echelle{libellé du 1}{libellé du 5}
\newcommand{\echelle}[2]{\par\nopagebreak\noindent$\square$~1\quad$\square$~2\quad$\square$~3\quad$\square$~4\quad$\square$~5\qquad{\small(1 = #1~; 5 = #2)}\par}

\begin{document}

\begin{center}
  {\Large\sffamily\bfseries Questionnaire~: communiquer sans réseau mobile}\\[4pt]
  {\sffamily Étude quantitative --- INSA Centre Val de Loire, STI 5A ASL-IA}
\end{center}

%==========================================================================
\section*{Notice pour l'équipe (à ne pas montrer aux répondants)}

\textbf{Objectif.} Mesurer, segment par segment, la chaîne \emph{problème $\rightarrow$ besoin $\rightarrow$
solution $\rightarrow$ intérêt $\rightarrow$ adoption $\rightarrow$ prix}, puis comparer les trois segments
(plein air~; camping, voyage, famille~; secours et gestion de crise).

\textbf{Principes de rédaction.}
\begin{itemize}[nosep]
  \item Les sections 1 à 4 viennent \emph{avant} la présentation du concept, pour ne pas influencer les réponses sur le problème et le besoin.
  \item Le concept est présenté sans nom, sans marque et sans prix~; le prix n'apparaît qu'en section 9, après la question ouverte (Q19).
  \item Les limites sont dites dans le concept (texte seul, pas une balise de détresse) pour éviter de sur-vendre.
  \item Une question = une idée~; échelles à cinq niveaux, ordre des cases tiré au hasard quand c'est possible~; toujours une case «~Autre~» ou «~Aucun~».
\end{itemize}

\textbf{Échantillon.} Non aléatoire, par choix raisonné. Tailles proposées~: 150 réponses pour le plein air,
100 pour camping, voyage et famille, 60 pour secours et gestion de crise (310 au total). La marge d'erreur d'un
pourcentage vaut environ $\pm\,98/\sqrt{n}$ points (95\,\%, tirage aléatoire supposé)~: $\pm\,8$ points pour $n=150$,
$\pm\,10$ pour $n=100$, $\pm\,13$ pour $n=60$. L'échantillon sert à \emph{comparer} les segments, pas à prévoir
les ventes ni à mesurer une part de marché.

\textbf{Où diffuser.} Plein air~: clubs de randonnée et de montagne, groupes de pratiquants en ligne, magasins outdoor,
refuges (affiche avec QR code). Camping, voyage, famille~: campings (QR code à l'accueil), groupes de camping-caristes
et de voyageurs, associations de parents, réseau personnel. Secours et gestion de crise~: ADRASEC et associations de
radioamateurs, Protection civile, Croix-Rouge, mairies, réseaux d'ONG. Les liens ou QR codes portent un paramètre
(\texttt{?src=club}, \texttt{?src=camping}\dots) pour compter les réponses par canal.

\textbf{Avant de diffuser.} Pré-test sur dix personnes (durée, questions mal comprises). Réponses anonymes~;
l'adresse e-mail de la question 23 est facultative et stockée à part (RGPD). Nettoyage à prévoir~: questionnaires
terminés en moins de trois minutes, prix maximal de la question 19 supérieur à 500\,\texteuro{} (à traiter à part).

\newpage
%==========================================================================
\section*{Questionnaire}

\begin{center}
  \fbox{\parbox{0.92\linewidth}{\small
  Bonjour, nous sommes des étudiants de l'INSA Centre Val de Loire et menons une enquête sur la communication en dehors
  des zones couvertes par le réseau mobile. Elle dure environ 8 à 10 minutes. Vos réponses sont anonymes, il n'y a pas
  de bonne ou de mauvaise réponse, et vous pouvez arrêter à tout moment. Cette enquête s'adresse aux personnes majeures.}}
\end{center}

%--------------------------------------------------------------------------
\sect{1}{Votre profil}

\quest{1}{Parmi ces profils, lequel vous correspond le mieux~?}
\logique{une seule réponse~; détermine le segment}
\begin{choix}
  \item Plein air et sport outdoor (randonnée, trail, VTT, alpinisme, ski, escalade\dots)
  \item Camping, voyage ou sorties en pleine nature, seul(e), entre amis ou en famille
  \item Secours et gestion de crise (professionnel, bénévole, association, commune, ONG)
  \item Autre (précisez)~: \rule{4cm}{0.3pt} \ \logique{analysé à part}
\end{choix}
\logique{profil 3 $\rightarrow$ poser les variantes de l'annexe A}

\quest{2}{Quelle est votre tranche d'âge~?}
\begin{choix}\setlength{\itemsep}{0pt}
  \item 18 à 24 ans \quad $\square$ 25 à 34 ans \quad $\square$ 35 à 44 ans \quad $\square$ 45 à 54 ans \quad $\square$ 55 à 64 ans \quad $\square$ 65 ans et plus
\end{choix}

\quest{3}{Où vivez-vous~?}
\begin{choix}
  \item En ville \quad $\square$ En zone périurbaine \quad $\square$ À la campagne \quad $\square$ En montagne
\end{choix}

%--------------------------------------------------------------------------
\sect{2}{Habitudes et pratiques}

\quest{4}{Au cours des 12 derniers mois, quelles activités avez-vous pratiquées en pleine nature~?}
\logique{plusieurs réponses~; si «~Aucune~» $\rightarrow$ passer à Q7}
\begin{choixdeux}
  \item Randonnée, trekking
  \item Trail, course à pied nature
  \item VTT, vélo
  \item Alpinisme, escalade
  \item Ski, raquettes
  \item Camping, bivouac, camping-car
  \item Voyage en zone isolée
  \item Chasse, pêche
  \item Autre (précisez)~: \rule{4cm}{0.3pt}
  \item Aucune
\end{choixdeux}

\quest{5}{En moyenne, à quelle fréquence pratiquez-vous ces activités~?}
\begin{choix}
  \item Quelques fois par an \quad $\square$ Une à trois fois par mois \quad $\square$ Chaque semaine \quad $\square$ Plusieurs fois par semaine
\end{choix}

\quest{6}{Le plus souvent, vous sortez\dots}
\begin{choix}
  \item seul(e) \quad $\square$ à deux \quad $\square$ en groupe de 3 à 5 personnes \quad $\square$ en groupe de 6 personnes ou plus \quad $\square$ cela dépend
\end{choix}

%--------------------------------------------------------------------------
\sect{3}{Existence du problème}

\quest{7}{Vous arrive-t-il de vous retrouver dans une zone sans réseau mobile (ni appel, ni SMS possible)~?}
\logique{si «~Jamais~» $\rightarrow$ passer à Q10}
\begin{choix}
  \item Jamais \quad $\square$ Rarement \quad $\square$ Parfois \quad $\square$ Souvent \quad $\square$ Très souvent
\end{choix}

\quest{8}{Où cela vous est-il arrivé~?}
\logique{plusieurs réponses}
\begin{choixdeux}
  \item En montagne
  \item À la campagne ou en forêt
  \item En mer ou sur le littoral
  \item À l'étranger
  \item En ville (sous-sol, bâtiment, foule)
  \item Lors d'une catastrophe ou d'une panne de réseau
  \item Autre (précisez)~: \rule{4cm}{0.3pt}
\end{choixdeux}

\quest{9}{Dans ces moments, comment communiquez-vous aujourd'hui avec votre groupe ou vos proches~?}
\logique{plusieurs réponses}
\begin{choixdeux}
  \item Je ne communique pas
  \item Des rendez-vous ou des horaires convenus à l'avance
  \item Des talkies-walkies
  \item La radio amateur
  \item Une balise de détresse
  \item Un communicateur satellite (messages par satellite)
  \item Le SOS satellite de mon smartphone
  \item Autre (précisez)~: \rule{4cm}{0.3pt}
\end{choixdeux}

%--------------------------------------------------------------------------
\sect{4}{Importance du besoin}

\quest{10}{Quand il n'y a pas de réseau, à quel point est-il important pour vous de pouvoir\dots}
\logique{ordre des lignes tiré au hasard}
\par\nopagebreak\smallskip
\begingroup\small\renewcommand{\arraystretch}{1.25}
\begin{tabular}{@{}p{8.2cm}ccccc@{}}
  & 1 & 2 & 3 & 4 & 5 \\
  \midrule
  a)~rester en contact avec les personnes qui sont avec vous & $\square$ & $\square$ & $\square$ & $\square$ & $\square$ \\
  b)~prévenir vos proches que tout va bien & $\square$ & $\square$ & $\square$ & $\square$ & $\square$ \\
  c)~partager votre position & $\square$ & $\square$ & $\square$ & $\square$ & $\square$ \\
  d)~coordonner une équipe (si cela vous concerne) & $\square$ & $\square$ & $\square$ & $\square$ & $\square$ \\
  e)~alerter les secours & $\square$ & $\square$ & $\square$ & $\square$ & $\square$ \\
\end{tabular}
\endgroup\par
{\small(1 = pas du tout important~; 5 = très important)}

\quest{11}{L'absence de réseau vous a-t-elle déjà posé un vrai problème (se perdre de vue, ne pas pouvoir prévenir quelqu'un, renoncer à une sortie\dots)~?}
\begin{choix}
  \item Jamais \quad $\square$ Une fois \quad $\square$ Plusieurs fois
\end{choix}
\logique{si «~Une fois~» ou «~Plusieurs fois~» $\rightarrow$ champ libre facultatif~: «~Décrivez en une phrase la dernière fois.~»}

%--------------------------------------------------------------------------
\vspace{10pt}
\begin{center}
  \fbox{\parbox{0.92\linewidth}{\small
  \textbf{Présentation d'une solution}\quad\logique{page séparée, avec une image sans marque}\\[3pt]
  Imaginez un petit appareil radio de poche, relié à votre téléphone par Bluetooth. Il permet d'envoyer des messages texte
  à des proches qui ont le même appareil, sans réseau mobile, sans internet et sans abonnement. Les messages passent d'un
  appareil à l'autre par radio~: la portée est de 0,5 à 2\,km en ville et de 5 à 15\,km en terrain dégagé, davantage si d'autres appareils relaient.
  Texte seulement~: pas d'appel, pas de photo. Ce n'est pas une balise de détresse~: l'appareil ne prévient pas les secours.\\[3pt]
  Les questions suivantes portent sur cette solution.}}
\end{center}

%--------------------------------------------------------------------------
\sect{5}{Intérêt pour la solution}

\quest{12}{Dans quelle mesure cette solution vous paraît-elle utile pour vous~?}
\echelle{pas du tout utile}{très utile}

\quest{13}{Qu'est-ce qui vous intéresse le plus~?}
\logique{deux réponses au plus~; ordre tiré au hasard}
\begin{choixdeux}
  \item Elle fonctionne sans réseau
  \item Il n'y a pas d'abonnement
  \item Mes messages restent privés (pas de compte, pas de serveur central)
  \item Elle est simple à utiliser
  \item Son prix
  \item Autre (précisez)~: \rule{4cm}{0.3pt}
  \item Rien ne m'intéresse
\end{choixdeux}

%--------------------------------------------------------------------------
\sect{6}{Usages envisagés}

\quest{14}{Dans quelles situations l'utiliseriez-vous~?}
\logique{plusieurs réponses}
\begin{choixdeux}
  \item Randonnée, montagne
  \item Camping, bivouac
  \item Voyage
  \item Sport (trail, VTT, ski\dots)
  \item Zones isolées, campagne
  \item Intervention de secours
  \item Catastrophe ou panne de réseau
  \item Événement ou festival
  \item Autre (précisez)~: \rule{4cm}{0.3pt}
  \item Je ne l'utiliserais pas
\end{choixdeux}

%--------------------------------------------------------------------------
\sect{7}{Freins à l'adoption}

\quest{15}{Qu'est-ce qui pourrait vous empêcher de l'acheter~?}
\logique{plusieurs réponses~; ordre tiré au hasard}
\begin{choixdeux}
  \item Le prix
  \item Devoir équiper aussi les autres personnes
  \item La portée
  \item L'autonomie de la batterie
  \item La complexité d'utilisation
  \item Elle n'est pas assez utile pour moi
  \item J'ai déjà une solution
  \item La confiance dans une marque inconnue
  \item Autre (précisez)~: \rule{4cm}{0.3pt}
  \item Rien
\end{choixdeux}

\quest{16}{Parmi ces freins, lequel est le principal~?}
\logique{un seul choix, parmi les réponses données à Q15}

%--------------------------------------------------------------------------
\sect{8}{Intention d'achat ou d'adoption}

\quest{17}{Seriez-vous prêt(e) à acheter ou à adopter cette solution~?}
\begin{choix}
  \item Certainement pas \quad $\square$ Probablement pas \quad $\square$ Peut-être \quad $\square$ Probablement \quad $\square$ Certainement
\end{choix}

\quest{18}{Combien de personnes de votre entourage équiperiez-vous, vous compris~?}
\begin{choix}
  \item 1 \quad $\square$ 2 \quad $\square$ 3 à 4 \quad $\square$ 5 à 9 \quad $\square$ 10 ou plus
\end{choix}

%--------------------------------------------------------------------------
\sect{9}{Prix acceptable}

\quest{19}{Quel prix maximal seriez-vous prêt(e) à payer pour un appareil, sans abonnement~?}
\logique{réponse libre, en euros~; à poser \emph{avant} toute mention d'un prix}\par
\noindent\rule{3cm}{0.3pt}~\texteuro

\quest{20}{Cet appareil coûterait 69\,\texteuro, en un seul achat et sans abonnement. L'achèteriez-vous~?}
\logique{variante possible~: tirer au hasard 59, 69 ou 79\,\texteuro{} pour lire trois points de la courbe de demande}
\begin{choix}
  \item Oui \quad $\square$ Peut-être \quad $\square$ Non
\end{choix}

\quest{21}{Pour équiper plusieurs personnes, quelle formule préféreriez-vous~?}
\logique{prix de l'équipe, à aligner sur le prix unitaire retenu~; si l'on craint un effet d'ancrage avec Q20, répartir au hasard (moitié Q20, moitié Q21)}
\begin{choix}
  \item Solo~: 1 appareil à 79\,\texteuro
  \item Duo~: 2 appareils à 139\,\texteuro
  \item Groupe~: 4 appareils à 249\,\texteuro
  \item Aucune de ces formules
\end{choix}

\quest{22}{Où l'achèteriez-vous le plus volontiers~?}
\begin{choixdeux}
  \item Dans une boutique de sport ou d'outdoor
  \item Sur le site internet de la marque
  \item Par mon club ou ma fédération (achat groupé)
  \item Par ma commune ou mon association (achat collectif)
  \item Sur un site de vente généraliste
  \item Autre (précisez)~: \rule{4cm}{0.3pt}
\end{choixdeux}

\quest{23}{(Facultatif) Souhaitez-vous être informé(e) du projet ou tester l'appareil gratuitement~?}
\begin{choix}
  \item Non, merci
  \item Oui, mon adresse e-mail~: \rule{6cm}{0.3pt}
\end{choix}
{\small Cette adresse sera utilisée uniquement pour vous recontacter au sujet de ce projet d'étude et supprimée à la fin de l'étude.}

\vspace{4pt}
\begin{center}
  \fbox{\parbox{0.92\linewidth}{\small Merci pour votre participation~! Vos réponses nous aident à savoir si cette solution répond à un vrai besoin.}}
\end{center}

\newpage
%==========================================================================
\section*{Annexe A~: variante « secours et gestion de crise » (profil 3)}

Mêmes sections, mêmes numéros. Seules les questions suivantes sont reformulées~; les autres sont posées telles quelles.

\begingroup\small\renewcommand{\arraystretch}{1.35}
\begin{tabularx}{\linewidth}{@{}p{0.9cm}>{\raggedright\arraybackslash}X>{\raggedright\arraybackslash}X@{}}
  \toprule
  \textbf{Q} & \textbf{Énoncé} & \textbf{Réponses} \\
  \midrule
  \qa{4} & Dans quel cadre intervenez-vous ou participez-vous à la gestion de crise~? (plusieurs réponses) &
      secours en montagne ou en milieu isolé~; sécurité civile, pompiers~; radioamateurs (réseau d'urgence)~; commune ou collectivité~; ONG, humanitaire~; association de secouristes~; autre \\
  \qa{5} & Combien de fois par an intervenez-vous ou participez-vous à un exercice~? &
      moins de 5~; 5 à 20~; plus de 20 \\
  \qa{6} & Quelle est la taille habituelle de votre équipe sur le terrain~? &
      1 à 2~; 3 à 5~; 6 à 10~; plus de 10 \\
  \qa{7} & Lors de vos interventions ou de vos exercices, le réseau mobile est-il absent, coupé ou saturé~? &
      jamais $\rightarrow$ très souvent \\
  \qa{8} & Dans quelles situations~? (plusieurs réponses) &
      zone isolée~; coupure d'électricité~; catastrophe naturelle~; saturation du réseau lors d'un rassemblement~; zone de conflit~; autre \\
  \qa{9} & Quels moyens utilisez-vous alors~? (plusieurs réponses) &
      radio professionnelle (réseau dédié)~; talkies-walkies~; radio amateur~; téléphone satellite~; messagerie du smartphone quand c'est possible~; messagers à pied ou en véhicule~; autre \\
  \qa{11} & Cette absence de communication a-t-elle déjà compliqué une intervention ou un exercice~? &
      jamais~; une fois~; plusieurs fois \\
  \qa{14} & Dans quelles situations votre structure l'utiliserait-elle~? &
      interventions~; exercices~; renfort des moyens existants~; catastrophe ou panne de réseau~; événement~; jamais \\
  \qa{17} & Votre structure ou votre équipe serait-elle prête à acheter ou à adopter cette solution~? &
      certainement pas $\rightarrow$ certainement \\
  \qa{18} & Combien d'appareils faudrait-il pour votre équipe~? &
      1 à 2~; 3 à 5~; 6 à 10~; 11 à 30~; plus de 30 \\
  \bottomrule
\end{tabularx}
\endgroup

\vspace{6pt}
Pour les structures qui achètent (communes, associations), ajouter si besoin~: «~Qui décide d'un achat de ce type, et avec quel budget~?~» (réponse libre).

\newpage
%==========================================================================
\section*{Annexe B~: indicateurs de lecture et seuils proposés}

Les seuils sont des propositions de départ, à valider \emph{avant} l'enquête~; ils servent à comparer les segments.
Un segment est dit prioritaire s'il atteint au moins six seuils sur huit~; si le «~oui~» à 69\,\texteuro{} reste sous 10\,\%
dans tous les segments, revoir le prix ou le concept avant le moule et la série.

\begingroup\small\renewcommand{\arraystretch}{1.3}
\begin{tabularx}{\linewidth}{@{}>{\raggedright\arraybackslash}X>{\raggedright\arraybackslash}p{4.3cm}>{\centering\arraybackslash}p{1.8cm}@{}}
  \toprule
  Indicateur & Calcul & Seuil proposé \\
  \midrule
  Problème~: sans réseau «~souvent~» ou «~très souvent~» & part des réponses «~Souvent~» et «~Très souvent~» à Q7 & $\geq$ 40\,\% \\
  Besoin~: «~important~» ou «~très important~» & part des notes 4 ou 5 à Q10\,a ou Q10\,b & $\geq$ 60\,\% \\
  Problème déjà vécu & part de «~Une fois~» et «~Plusieurs fois~» à Q11 & $\geq$ 30\,\% \\
  Utilité & part des notes 4 ou 5 à Q12 & $\geq$ 50\,\% \\
  Intention d'achat & part de «~Probablement~» et «~Certainement~» à Q17 & $\geq$ 30\,\% \\
  Prix~: «~oui~» à 69\,\texteuro & part de «~Oui~» à Q20 & $\geq$ 20\,\% \\
  Prix maximal & médiane de Q19 (valeurs $>$ 500\,\texteuro{} écartées) & $\geq$ 60\,\texteuro \\
  Appareils à équiper & médiane de Q18 (classes codées 1, 2, 3,5, 7, 10) & $\geq$ 2 \\
  \bottomrule
\end{tabularx}
\endgroup

\vspace{6pt}
\textbf{Lectures complémentaires.}
\begin{itemize}[nosep]
  \item \emph{Part solvable à 69\,\texteuro}~: part des réponses de Q19 au moins égales à 69\,\texteuro.
  \item \emph{Recette à un prix $p$}~: $p \times$ (part des réponses de Q19 au moins égales à $p$)~; comparer 59, 69 et 79\,\texteuro.
  \item \emph{Effet du prix}~: part de «~Probablement~» et «~Certainement~» à Q17, moins part de «~Oui~» à Q20.
  \item \emph{Formule préférée}~: répartition de Q21 (Solo, Duo, Groupe, aucune), croisée avec Q18.
  \item \emph{Canaux}~: répartition de Q22, par segment, pour choisir où vendre en premier.
\end{itemize}

\end{document}
